\documentclass[12pt,a4paper]{amsart}
\usepackage{amsmath,amssymb,amsthm}
\usepackage{mathtools}
\usepackage[utf8]{inputenc}
\usepackage[T1]{fontenc}
\usepackage{hyperref}
\usepackage{enumitem,booktabs,array}
\usepackage{geometry}
\geometry{margin=2.5cm}

\newtheorem{theorem}{Satz}[section]
\newtheorem{lemma}[theorem]{Lemma}
\newtheorem{corollary}[theorem]{Korollar}
\newtheorem{proposition}[theorem]{Proposition}
\theoremstyle{definition}
\newtheorem{definition}[theorem]{Definition}
\newtheorem{example}[theorem]{Beispiel}
\newtheorem{remark}[theorem]{Bemerkung}

\newtheorem{conjecture}[theorem]{Vermutung}

\author{Michael Fuhrmann}
\address{Specialist Mathematics Project, Build 133}
\email{specialist-maths@localhost}
\date{\today}

\title{Die Lehmer-Vermutung über die Ramanujan-$\tau$-Funktion:\\
Modularität, Nicht-Verschwinden und numerische Evidenz}

\begin{document}

\begin{abstract}
Die Lehmer-Vermutung, aufgestellt 1947 von Derrick Henry Lehmer, behauptet, dass
die Ramanujan-$\tau$-Funktion die Eigenschaft $\tau(n) \neq 0$ für alle positiven
ganzen Zahlen $n$ besitzt. Trotz nahezu achtzig Jahren intensiver Forschung ist
diese verblüffend einfach formulierte Aussage bis heute unbewiesen. In dieser Arbeit
entwickeln wir den theoretischen Rahmen rund um die Vermutung: die Arithmetik der
Diskriminantenmodulform $\Delta$, die Hecke-Theorie, die Multiplikativität von
$\tau$, das Eulerprodukt von $L(s, \Delta)$ und Delignes berühmten Satz
$|\tau(p)| \leq 2p^{11/2}$, der die Ramanujan-Vermutung über die Größenordnung von
$\tau$ bewiesen hat. Wir präsentieren modulare Kongruenzen, strukturelle Konsequenzen
einer hypothetischen Nullstelle sowie eine numerische Verifikation im Specialist
Mathematics Project bis $n \leq 50{,}000$, die $\tau(n) \neq 0$ in diesem Bereich
bestätigt. Der Zusammenhang zwischen der Lehmer-Vermutung und mod-$\ell$
Galois-Darstellungen nach Serre und Swinnerton-Dyer wird skizziert. Wir schließen,
dass die Vermutung zwar durch umfangreiche Rechnungen gut gestützt ist, aber eines
der tiefsten offenen Probleme in der Theorie der Modulformen bleibt.
\end{abstract}

\maketitle

\tableofcontents

%% ============================================================
\section{Einleitung}
%% ============================================================

In seiner Arbeit \emph{On certain arithmetical functions} von 1916 führte Srinivasa
Ramanujan die arithmetische Funktion $\tau : \mathbb{N} \to \mathbb{Z}$ durch die
erzeugende Reihe der Diskriminantenmodulform ein:
\begin{equation}\label{eq:delta-def}
  \Delta(q) \;=\; q \prod_{n=1}^{\infty}(1-q^n)^{24}
  \;=\; \sum_{n=1}^{\infty} \tau(n)\, q^n,
  \qquad |q| < 1.
\end{equation}
Ramanujan berechnete die ersten Werte von $\tau$ von Hand und entdeckte dabei
mehrere bemerkenswerte arithmetische Eigenschaften: Multiplikativität, Kongruenzen
modulo kleiner Primzahlen und --- auffallend --- dass kein einziger der von ihm
berechneten Werte null war.

Dreißig Jahre später nahm Derrick Henry Lehmer \cite{Lehmer1947} diese Berechnungen
wieder auf und formulierte folgende Vermutung:

\begin{center}
  \textit{``Es liegt nahe zu fragen, ob $\tau(n)$ jemals null sein kann. Es wurde
  kein Wert von $n$ gefunden, für den $\tau(n) = 0$ gilt, und es erscheint sehr
  unwahrscheinlich, dass ein solches $n$ existiert.''}
\end{center}

Diese Aussage --- heute als \emph{Lehmer-Vermutung} bekannt --- hat sich seit mehr
als fünfundsiebzig Jahren jedem Beweisversuch widersetzt. Sie steht an der
Schnittstelle von analytischer Zahlentheorie, Theorie der Modulformen und
arithmetischer Geometrie und bedient sich Werkzeugen von Hecke-Operatoren bis hin
zu Galois-Darstellungen.

Ziel dieser Arbeit ist dreifach. Erstens entwickeln wir den Modulformen-Hintergrund,
der zum Verständnis der Vermutung notwendig ist. Zweitens geben wir einen Überblick
über die wichtigsten theoretischen Ergebnisse: Delignes Beweis der
Ramanujan-Vermutung, die Hecke-Eigenwert-Struktur von $\Delta$ und die Verbindungen
zu mod-$\ell$ Galois-Darstellungen nach Serre und Swinnerton-Dyer. Drittens
präsentieren wir numerische Evidenz aus dem Specialist Mathematics Project: eine
direkte Verifikation, dass $\tau(n) \neq 0$ für alle $n \leq 50{,}000$ gilt,
durchgeführt mittels exakter BigInt-Polynomialentwicklung von $\Delta(q)$.

\medskip
\noindent\textbf{Notation.} Im gesamten Text bezeichnet $\mathbb{H} = \{z \in
\mathbb{C} : \Im z > 0\}$ die obere Halbebene, $q = e^{2\pi i z}$ für
$z \in \mathbb{H}$, und $\mathrm{SL}_2(\mathbb{Z})$ ist die volle Modulgruppe.
Für eine ganze Zahl $k$ und eine positive ganze Zahl $n$ sei
$\sigma_k(n) = \sum_{d \mid n} d^k$.

%% ============================================================
\section{Hintergrund: Modulformen}
%% ============================================================

\subsection{Die obere Halbebene und die Modulgruppe}

Die Modulgruppe $\Gamma = \mathrm{SL}_2(\mathbb{Z})$ operiert auf $\mathbb{H}$ durch
Möbius-Transformationen:
\[
  \begin{pmatrix} a & b \\ c & d \end{pmatrix} \cdot z
  \;=\; \frac{az+b}{cz+d}.
\]
Die Erzeuger sind $S: z \mapsto -1/z$ und $T: z \mapsto z+1$, mit Fundamentalbereich
$\mathcal{F} = \{z \in \mathbb{H} : |\Re z| \leq 1/2,\ |z| \geq 1\}$.

\begin{definition}
Eine \emph{Modulform vom Gewicht $k$ für $\Gamma$} ist eine holomorphe Funktion
$f : \mathbb{H} \to \mathbb{C}$, die die Transformationseigenschaft
\[
  f\!\left(\frac{az+b}{cz+d}\right) = (cz+d)^k f(z)
  \quad \text{für alle } \begin{pmatrix}a&b\\c&d\end{pmatrix} \in \Gamma
\]
erfüllt und holomorph in der Spitze $i\infty$ ist. Verschwindet $f$ zusätzlich in
der Spitze, so heißt sie \emph{Spitzenform}. Den Raum der Spitzenformen vom Gewicht
$k$ für $\Gamma$ bezeichnen wir mit $S_k(\Gamma)$.
\end{definition}

Die Dimensionsformel liefert $\dim S_{12}(\Gamma) = 1$, wodurch $\Delta$ eine
einzigartige Stellung innehat: es ist (bis auf Skalierung) die \emph{einzige}
normalisierte Spitzenform vom Gewicht $12$.

\subsection{Eisenstein-Reihen und der Ring der Modulformen}

Der graduierte Ring $M_*(\Gamma) = \bigoplus_{k \geq 0} M_k(\Gamma)$ wird frei von
den Eisenstein-Reihen
\[
  E_4(z) = 1 + 240\sum_{n=1}^\infty \sigma_3(n)q^n, \qquad
  E_6(z) = 1 - 504\sum_{n=1}^\infty \sigma_5(n)q^n
\]
erzeugt. Jede Modulform für $\Gamma$ ist ein Polynom in $E_4$ und $E_6$.
Die Diskriminantenform lässt sich schreiben als
\begin{equation}\label{eq:delta-E4E6}
  \Delta(z) = \frac{E_4(z)^3 - E_6(z)^2}{1728}.
\end{equation}
Diese Identität, zusammen mit den Verschwindungsordnungen von $E_4$ und $E_6$ in
elliptischen Punkten, zeigt, dass $\Delta$ keine Nullstellen in $\mathbb{H}$ hat;
seine einzige Nullstelle auf der kompaktifizierten Modulkurve $X(1)$ liegt in der
Spitze.

\subsection{Die Dedekind-Eta-Funktion}

Die Dedekind-Eta-Funktion
\[
  \eta(z) = e^{\pi i z/12} \prod_{n=1}^\infty (1 - e^{2\pi i n z})
  = q^{1/24} \prod_{n=1}^\infty (1 - q^n)
\]
erfüllt $\Delta(z) = (2\pi)^{12} \eta(z)^{24}$ und transformiert sich unter der
Modulgruppe gemäß
\[
  \eta\!\left(\frac{az+b}{cz+d}\right)
  = \varepsilon(a,b,c,d)\,(cz+d)^{1/2}\,\eta(z),
\]
wobei $\varepsilon$ eine $24$-ste Einheitswurzel ist. Die Produktformel impliziert
sofort $\eta(z) \neq 0$ für alle $z \in \mathbb{H}$. Dies ist eines der wenigen
Nicht-Verschwindenresultate, die wir für $\Delta$ haben --- jedoch sagt es nichts
über einzelne Fourier-Koeffizienten $\tau(n)$ aus.

%% ============================================================
\section{Die Ramanujan-$\tau$-Funktion}
%% ============================================================

\subsection{Definition und erste Werte}

Die Ramanujan-$\tau$-Funktion ist durch die Fourier-Entwicklung \eqref{eq:delta-def}
definiert. Die ersten Werte, berechnet durch direkte Produktentwicklung, sind in
Tabelle~\ref{tab:tau-werte} aufgelistet.

\begin{table}[ht]
\centering
\caption{Erste Werte der Ramanujan-$\tau$-Funktion}
\label{tab:tau-werte}
\begin{tabular}{@{}rr@{\quad}rr@{}}
\toprule
$n$ & $\tau(n)$ & $n$ & $\tau(n)$ \\
\midrule
1  & $1$           & 7  & $-16{.}744$   \\
2  & $-24$         & 8  & $84{.}480$    \\
3  & $252$         & 9  & $-113{.}643$  \\
4  & $-1{.}472$    & 10 & $-115{.}920$  \\
5  & $4{.}830$     & 11 & $534{.}612$   \\
6  & $-6{.}048$    & 12 & $-370{.}944$  \\
\bottomrule
\end{tabular}
\end{table}

\noindent Weitere verifizierte Werte: $\tau(100) = 37{.}534{.}859{.}200$ und
$\tau(1000) = -30{.}328{.}412{.}970{.}240{.}000$.

\subsection{Multiplikativität}

Ramanujan vermutete, und Mordell \cite{Mordell1917} bewies 1917, dass $\tau$ eine
\emph{multiplikative} Funktion ist.

\begin{theorem}[Mordell, 1917]\label{thm:multiplikativitaet}
  Die Funktion $\tau$ ist multiplikativ: Wenn $\gcd(m,n) = 1$, dann gilt
  \[
    \tau(mn) = \tau(m)\,\tau(n).
  \]
  Für jede Primzahl $p$ und jedes $k \geq 1$ gilt außerdem die Rekurrenz
  \begin{equation}\label{eq:hecke-rekurrenz}
    \tau(p^{k+1}) = \tau(p)\,\tau(p^k) - p^{11}\,\tau(p^{k-1}).
  \end{equation}
\end{theorem}

\begin{proof}[Beweisidee]
  Mordells ursprüngliches Argument benutzt die Theorie der Hecke-Operatoren; siehe
  Abschnitt~\ref{sec:hecke} für Details. Die Rekurrenz \eqref{eq:hecke-rekurrenz}
  folgt aus der multiplikativen Struktur der Hecke-Eigenwerte.
\end{proof}

Eine grundlegende Konsequenz der Multiplikativität: $\tau(n) = 0$ für \emph{irgend}ein
$n$ genau dann, wenn $\tau(p^k) = 0$ für eine Primzahlpotenz $p^k$. Durch
\eqref{eq:hecke-rekurrenz} wird $\tau(p^k) = 0$ vollständig durch $\tau(p)$
bestimmt. Dies reduziert die Lehmer-Vermutung auf:
\begin{equation}\label{eq:reduzierte-vermutung}
  \tau(p) \neq 0 \quad \text{für alle Primzahlen } p.
\end{equation}

\subsection{Wachstum und die Ramanujan-Vermutung}

Ramanujan beobachtete empirisch, dass $|\tau(p)|$ wie $p^{11/2}$ wächst, und
vermutete $|\tau(p)| \leq 2p^{11/2}$ für alle Primzahlen $p$. Diese Schranke ---
alles andere als offensichtlich --- wurde von Deligne als Konsequenz der
Weil-Vermutungen bewiesen; Abschnitt~\ref{sec:deligne} ist diesem Ergebnis
gewidmet.

%% ============================================================
\section{Die Lehmer-Vermutung}
%% ============================================================

\subsection{Formale Aussage}

\begin{definition}
  Die \emph{Ramanujan-$\tau$-Funktion} ist die arithmetische Funktion
  $\tau : \mathbb{N} \to \mathbb{Z}$, definiert durch \eqref{eq:delta-def}.
\end{definition}

\begin{remark}
  Da $\Delta$ ganzzahlige Fourier-Koeffizienten hat und $\Delta(z) \neq 0$ für alle
  $z \in \mathbb{H}$ gilt, sind die Werte $\tau(n)$ wohldefinierte, nichtverschwindende
  algebraische Objekte im funktionentheoretischen Sinne. Die Vermutung betrifft
  jedoch die arithmetischen Werte der Fourier-Koeffizienten.
\end{remark}

\begin{conjecture}[Lehmer, 1947]\label{conj:lehmer}
  Für alle $n \geq 1$ gilt
  \[
    \tau(n) \neq 0.
  \]
\end{conjecture}

Dies ist eine Aussage über ganzzahlige Werte einer spezifischen arithmetischen
Funktion; sie folgt weder aus dem funktionentheoretischen Nicht-Verschwinden von
$\Delta$ auf $\mathbb{H}$ noch aus irgendeiner derzeit bekannten analytischen
Eigenschaft von $L(s, \Delta)$.

\subsection{Historischer Hintergrund}

Lehmer \cite{Lehmer1947} selbst verifizierte die Vermutung für
$n \leq 3{.}316{.}799$ und bemerkte die Schwierigkeit eines Beweises. Seitdem hat
sich die Verifikationsgrenze stetig weiterentwickelt:
\begin{itemize}[leftmargin=2em]
  \item Serre \cite{Serre1987} stellte fest, dass eine Nullstelle $\tau(p) = 0$
        für eine Primzahl $p$ drastische Konsequenzen für die mod-$\ell$
        Galois-Darstellungen von $\Delta$ hätte.
  \item Swinnerton-Dyer \cite{SwinnertonDyer1973} untersuchte die Kongruenzen von
        $\tau$ modulo verschiedener Primzahlen und ihre Implikationen.
  \item Chua \cite{Chua2004} erweiterte die Verifikation auf $n \approx 10^{23}$
        mithilfe von Hecke-Eigenwert-Techniken.
\end{itemize}

Die Vermutung bleibt offen (Stand 2026).

\subsection{Warum ist sie schwer?}

Die Schwierigkeit rührt daher, dass wir --- weder aus der algebraischen noch aus
der analytischen Theorie von $\Delta$ --- keinen strukturellen Grund kennen, warum
ein Koeffizient verschwinden sollte. Das funktionentheoretische Nicht-Verschwinden
von $\Delta(z)$ zeigt, dass das \emph{Produkt} aller Koeffizienten (kodiert in der
erzeugenden Funktion) ungleich null ist, aber einzelne Koeffizienten sind eine
ganz andere Sache.

Vergleiche dazu das Nicht-Verschwinden von Hecke-$L$-Funktionen bei $s=1$ (bewiesen
über die analytische Klassenzahlformel) oder das Nicht-Verschwinden von $\zeta(s)$
auf $\Re s = 1$ (Hadamard und de la Vallée Poussin): jene Beweise nutzen analytische
Eigenschaften der zugehörigen $L$-Funktion direkt. Für die Lehmer-Vermutung ist kein
analoger Ansatz bekannt.

%% ============================================================
\section{Hecke-Operatoren und die $L$-Funktion von $\Delta$}
\label{sec:hecke}
%% ============================================================

\subsection{Hecke-Operatoren}

Für eine positive ganze Zahl $n$ ist der \emph{Hecke-Operator} $T_n$ auf
$M_k(\Gamma)$ definiert durch
\[
  (T_n f)(z) = n^{k-1} \sum_{ad=n,\, b=0}^{d-1}
  d^{-k} f\!\left(\frac{az+b}{d}\right).
\]
Die Operatoren $\{T_n\}$ bilden einen kommutativen Ring, und $\Delta$ ist eine
Hecke-Eigenform: $T_n \Delta = \tau(n)\, \Delta$ für alle $n \geq 1$. Dies ist die
Quelle der Multiplikativität.

\begin{theorem}[Hecke]\label{thm:hecke-eigenform}
  $\Delta$ ist eine normalisierte Hecke-Eigenform in $S_{12}(\Gamma)$, d.\,h.
  \[
    T_n \Delta = \tau(n)\, \Delta \quad \text{für alle } n \geq 1.
  \]
\end{theorem}

\subsection{Das Eulerprodukt}

Die Multiplikativität der Hecke-Eigenwerte überträgt sich direkt in ein Eulerprodukt
für die zugehörige $L$-Funktion.

\begin{definition}
  Die \emph{$L$-Funktion von $\Delta$} ist
  \[
    L(s, \Delta) = \sum_{n=1}^{\infty} \frac{\tau(n)}{n^s},
    \qquad \Re s > 13/2 \text{ (absolute Konvergenz nach Deligne).}
  \]
\end{definition}

\begin{theorem}[Eulerprodukt]\label{thm:eulerprodukt}
  Für hinreichend großes $\Re s$ gilt
  \begin{equation}\label{eq:eulerprodukt}
    L(s, \Delta)
    = \prod_{p \text{ prim}}
      \frac{1}{1 - \tau(p)\,p^{-s} + p^{11-2s}}.
  \end{equation}
\end{theorem}

\begin{proof}
  Dies folgt aus der Multiplikativität von $\tau$ und der Rekurrenz
  \eqref{eq:hecke-rekurrenz}: Für jede Primzahl $p$ summiert man die geometrische
  Reihe $\sum_{k=0}^\infty \tau(p^k) X^k$ (mit $X = p^{-s}$) und erhält mithilfe
  der Rekurrenz den quadratischen Faktor $1 - \tau(p)X + p^{11}X^2$.
\end{proof}

\begin{remark}
  Jeder lokale Faktor $1 - \tau(p)p^{-s} + p^{11-2s}$ faktorisiert als
  $(1 - \alpha_p p^{-s})(1 - \beta_p p^{-s})$, wobei $\alpha_p \beta_p = p^{11}$
  und $\alpha_p + \beta_p = \tau(p)$. Delignes Satz besagt
  $|\alpha_p| = |\beta_p| = p^{11/2}$.
\end{remark}

\subsection{Konsequenz von $\tau(p_0) = 0$}

Wenn $\tau(p_0) = 0$ für eine Primzahl $p_0$, degeneriert der lokale Eulerprodukt-Faktor:
\[
  \frac{1}{1 - \tau(p_0)p_0^{-s} + p_0^{11-2s}}
  = \frac{1}{1 + p_0^{11-2s}}.
\]
Die Wurzeln $\alpha_{p_0} = ip_0^{11/2}$ und $\beta_{p_0} = -ip_0^{11/2}$ sind rein
imaginär --- ein strukturell höchst ungewöhnliches Verhalten im Vergleich zu allen
bisher bekannten Hecke-Eigenwerten.

%% ============================================================
\section{Delignes Satz}
\label{sec:deligne}
%% ============================================================

\subsection{Die Ramanujan-Vermutung}

Im Jahr 1916 vermutete Ramanujan aufgrund numerischer Evidenz, dass
$|\tau(p)| \leq 2p^{11/2}$ für alle Primzahlen $p$ gilt. Dies ist äquivalent dazu,
dass die Wurzeln $\alpha_p, \beta_p$ des Polynoms $X^2 - \tau(p)X + p^{11}$ auf dem
Kreis $|X| = p^{11/2}$ in $\mathbb{C}$ liegen.

\begin{theorem}[Deligne, 1974]\label{thm:deligne}
  Für jede Primzahl $p$ gilt
  \[
    |\tau(p)| \leq 2p^{11/2}.
  \]
\end{theorem}

\begin{proof}[Beweisüberblick]
  Deligne \cite{Deligne1974} bewies dies als Konsequenz seines Beweises der
  Weil-Vermutungen für glatte projektive Varietäten über endlichen Körpern.
  Die zentralen Schritte sind:
  \begin{enumerate}[label=(\roman*)]
    \item Man ordnet $\Delta$ eine $\ell$-adische Galois-Darstellung
          $\rho_\ell : \mathrm{Gal}(\overline{\mathbb{Q}}/\mathbb{Q}) \to
          \mathrm{GL}_2(\mathbb{Z}_\ell)$ zu.
    \item Es gilt $\tau(p) = \mathrm{Sp}(\rho_\ell(\mathrm{Frob}_p))$ für
          Primzahlen $p \nmid \ell$.
    \item Man zeigt, dass $\rho_\ell$ aus der Kohomologie einer Familie abelscher
          Varietäten stammt, sodass die Frobenius-Eigenwerte den Absolutwert
          $p^{11/2}$ haben (Riemann-Hypothese für abelsche Varietäten über
          $\mathbb{F}_p$, Weil-Vermutungen, bewiesen von Deligne in
          \cite{DelignePureI}).
  \end{enumerate}
\end{proof}

Deligne erhielt 1978 die Fields-Medaille, teilweise für diese Arbeit.

\begin{corollary}\label{cor:deligne-koeff}
  Für alle $n \geq 1$ gilt $|\tau(n)| \leq d(n)\, n^{11/2}$, wobei $d(n)$ die
  Anzahl der Teiler von $n$ ist.
\end{corollary}

\begin{remark}
  Delignes Satz zeigt, dass $\tau(p)$ nur dann nahe bei $0$ sein kann, wenn
  $\tau(p) \ll p^{11/2}$; er schließt jedoch $\tau(p) = 0$ \emph{exakt} nicht aus.
  Die Lehmer-Vermutung ist daher \emph{streng stärker} als die Ramanujan-Vermutung
  (die heute ein Satz ist).
\end{remark}

\subsection{Bezug zur Funktionalgleichung}

Die vervollständigte $L$-Funktion
\[
  \Lambda(s, \Delta) = \left(\frac{2\pi}{\sqrt{N}}\right)^{-s}
  \Gamma(s)\, L(s, \Delta), \qquad N = 1 \text{ für } \Delta,
\]
erfüllt die Funktionalgleichung
\[
  \Lambda(s, \Delta) = \Lambda(12-s, \Delta),
\]
die $s$ und $12-s$ symmetrisch um den Zentralpunkt $s = 6$ verbindet. Das
Nicht-Verschwinden von $L(6, \Delta)$ (und allgemeiner von $L(s, \Delta)$ auf der
kritischen Geraden $\Re s = 6$) ist ein separates, ebenfalls offenes Problem.

%% ============================================================
\section{Modulare Kongruenzen für $\tau$}
%% ============================================================

\subsection{Die Kongruenz modulo 691}

Die Primzahl $691$ spielt eine besondere Rolle in der Arithmetik der Modulformen.
Sie erscheint im Zähler der Bernoulli-Zahl $B_{12} = -691/2730$ und im konstanten
Term der Eisenstein-Reihe $E_{12}$.

\begin{theorem}[Ramanujan, 1916]\label{thm:kong-691}
  Für alle $n \geq 1$ gilt
  \[
    \tau(n) \equiv \sigma_{11}(n) \pmod{691}.
  \]
\end{theorem}

\begin{proof}[Beweisidee]
  Die einzige Spitzenform $\Delta \in S_{12}(\Gamma)$ und die Eisenstein-Reihe
  $E_{12}$ liegen beide in $M_{12}(\Gamma)$. Da $\dim M_{12}(\Gamma) = 2$ und der
  konstante Term von $E_{12}$ die Zahl $B_{12} = -691/2730$ enthält, teilt $691$
  die Differenz ihrer $q$-Entwicklungen koeffizientenweise.
\end{proof}

Diese Kongruenz war die historische Motivation für Serres Theorie der mod-$\ell$
Modulformen und Ribets Arbeit über Galois-Darstellungen, die schließlich zum Beweis
des Großen Fermatschen Satzes führte.

\subsection{Weitere Kongruenzen}

Über $691$ hinaus gibt es Kongruenzen für $\tau$ modulo kleiner Primzahlen:

\begin{proposition}\label{prop:kleine-prim-kong}
  Für alle $n \geq 1$ gelten die folgenden Kongruenzen:
  \begin{align}
    \tau(n) &\equiv n^2\,\sigma_7(n) \pmod{5}, \label{eq:kong5}\\
    \tau(n) &\equiv n\,\sigma_9(n) \pmod{7}. \label{eq:kong7}
  \end{align}
\end{proposition}

\begin{remark}
  Die rechten Seiten von \eqref{eq:kong5} und \eqref{eq:kong7} sind für $n$ mit
  $\gcd(n, 5) = 1$ bzw.\ $\gcd(n, 7) = 1$ ungleich null. Dennoch reichen die
  Kongruenzen allein nicht aus, um $\tau(n) \neq 0$ zu beweisen: Eine Kongruenz
  $\tau(n) \equiv r \pmod{\ell}$ mit $r \neq 0$ schließt $\tau(n) = 0$ nur für
  das jeweilige $n$ und $\ell$ aus. Man bräuchte unendlich viele solche Kongruenzen,
  um alle Fälle abzudecken.
\end{remark}

\subsection{Kongruenzen und Galois-Darstellungen}

Serre \cite{Serre1987} beobachtete, dass wenn $\tau(p_0) = 0$ für eine Primzahl
$p_0$ gilt, dann wäre die mod-$\ell$ Galois-Darstellung $\bar\rho_\ell$ von
$\Delta$ für $\ell = p_0$ reduzibel oder anderweitig degeneriert. Genauer:

\begin{proposition}[Serre]\label{prop:serre-galois}
  Angenommen, $\tau(p_0) = 0$ für eine Primzahl $p_0$. Sei $\ell \neq p_0$ eine
  Primzahl. Dann muss das charakteristische Polynom von $\mathrm{Frob}_{p_0}$
  auf $V_\ell(\Delta)$ von der Form $X^2 + p_0^{11}$ sein, d.\,h.\ $-p_0^{11}$
  muss modulo $\ell$ ein Quadrat sein --- für \emph{alle} Primzahlen $\ell$. Das
  ist eine außerordentlich starke Bedingung.
\end{proposition}

Dies liefert ein theoretisches Hindernis: Wenn $\tau(p_0) = 0$, muss $p_0$
gleichzeitig Kongruenzbedingungen modulo jeder Primzahl $\ell$ erfüllen, was im
Widerspruch zu Dirichlets Satz über Primzahlen in arithmetischen Progressionen
und Gleichverteilungsresultaten steht. Ein absoluter Widerspruch ist bisher jedoch
nicht hergeleitet worden.

%% ============================================================
\section{Numerische Verifikation}
%% ============================================================

\subsection{Algorithmus}

Wir beschreiben die numerische Verifikation, die im Specialist Mathematics Project
(Build 131, Skript \texttt{heavy\_compute\_lehmer\_tau.py}) durchgeführt wurde.

Der Kernalgorithmus ist die BigInt-Polynomialentwicklung:
\begin{enumerate}[label=\arabic*., leftmargin=2em]
  \item Darstellung von $\Delta(q) = q \prod_{k=1}^{N}(1-q^k)^{24}$ als formale
        Potenzreihe, abgebrochen bei Grad $N$.
  \item Entwicklung von $(1-q^k)^{24} = \sum_{j=0}^{24}\binom{24}{j}(-1)^j q^{kj}$
        für jedes $k = 1, \ldots, N$.
  \item Sequentielle Multiplikation dieser $N$ Polynome, unter Beibehaltung exakter
        BigInt-Koeffizienten (beliebig genaue ganze Zahlen).
  \item Extraktion von $\tau(n)$ als Koeffizient von $q^n$ im Ergebnispolynom.
\end{enumerate}

Es wird durchgehend exakte Ganzzahlarithmetik verwendet, sodass kein
Gleitkommafehler entstehen kann; jeder Wert $\tau(n)$ wird mit Sicherheit bestimmt.

\subsection{Ergebnisse}

\begin{theorem}[Numerisch, Build 131]\label{thm:numerisch}
  Folgendes wurde verifiziert:
  \begin{enumerate}[label=(\roman*)]
    \item $\tau(n) \neq 0$ für alle $n$ mit $1 \leq n \leq 50{.}000$.
    \item Die bekannten Werte aus Tabelle~\ref{tab:tau-werte}, einschließlich
          $\tau(100) = 37{.}534{.}859{.}200$ und
          $\tau(1000) = -30{.}328{.}412{.}970{.}240{.}000$, wurden bestätigt.
    \item Delignes Schranke $|\tau(p)| \leq 2p^{11/2}$ gilt für alle Primzahlen
          $p \leq 1{.}000$ (null Verletzungen gefunden).
  \end{enumerate}
  Gesamte Rechenzeit: $4{.}330$ Sekunden.
\end{theorem}

\begin{remark}
  Obwohl $n \leq 50{.}000$ bescheiden ist im Vergleich zum Stand der Technik
  (Chua \cite{Chua2004} verifizierte bis $n \approx 10^{23}$), dient die
  unabhängige Berechnung als Korrektheitsprüfung für den zugrunde liegenden
  Modulformen-Code im Specialist Mathematics Project und bestätigt die
  Implementierung von $\Delta$.
\end{remark}

\subsection{Stand der Technik}

Die umfassendste numerische Verifikation nutzt den Hecke-Eigenwert-Ansatz: Wenn
$\tau(n) = 0$ für irgendein $n$, dann insbesondere $\tau(p) = 0$ für eine Primzahl
$p \mid n$. Die Berechnung von Hecke-Eigenwerten --- statt der Fourier-Koeffizienten
direkt --- ermöglicht die Verifikation bis zu erheblich größeren Schranken. Chua
\cite{Chua2004} etablierte $\tau(n) \neq 0$ für
$n \leq 982{.}149{.}821{.}766{.}199{.}295{.}999 \approx 10^{23}$.
Dies bleibt der veröffentlichte Rekord (Stand 2026).

%% ============================================================
\section{Konsequenzen einer hypothetischen Nullstelle}
%% ============================================================

\subsection{Auswirkung auf das Eulerprodukt}

Angenommen, $\tau(n_0) = 0$ für irgendein $n_0$. Durch Multiplikativität und die
Diskussion in Abschnitt~\ref{sec:hecke} existiert eine Primzahl $p_0$ mit
$\tau(p_0) = 0$. Der lokale Eulerfaktor bei $p_0$ degeneriert zu
\[
  \frac{1}{1 + p_0^{11-2s}},
\]
welches Pole bei $s = 11/2 \pm \pi i / (2 \log p_0)$ hat. Das bedeutet:
$L(s, \Delta)$ hätte Eulerfaktornullstellen auf der kritischen Geraden
$\Re s = 11/2$. Zusammen mit der Funktionalgleichung würde dies starke
Einschränkungen für die globale Nullstellenverteilung von $L(s, \Delta)$ erzwingen.

\subsection{Auswirkung auf Galois-Darstellungen}

Die $\ell$-adische Darstellung $\rho_\ell : G_\mathbb{Q} \to \mathrm{GL}_2(\mathbb{Z}_\ell)$
von $\Delta$ erfüllt
\[
  \mathrm{Sp}(\rho_\ell(\mathrm{Frob}_p)) = \tau(p), \quad
  \det(\rho_\ell(\mathrm{Frob}_p)) = p^{11},
  \quad p \neq \ell.
\]
Wenn $\tau(p_0) = 0$ für $p_0 \neq \ell$, hat $\mathrm{Frob}_{p_0}$ die
Eigenwerte $\pm i p_0^{11/2}$ --- rein imaginär. Dies ist konsistent mit einer
irreduziblen $2$-dimensionalen Darstellung, aber das Verhalten wie eine
$90^\circ$-Drehung in $\mathrm{GL}_2$ ist äußerst nicht-generisch.

\subsection{Auswirkung auf höhere Koeffizienten}

Mithilfe der Rekurrenz \eqref{eq:hecke-rekurrenz} ergibt sich bei $\tau(p_0) = 0$:
\begin{align*}
  \tau(p_0^0) &= 1, \\
  \tau(p_0^1) &= 0, \\
  \tau(p_0^2) &= 0 \cdot 0 - p_0^{11} \cdot 1 = -p_0^{11} \neq 0, \\
  \tau(p_0^3) &= 0 \cdot (-p_0^{11}) - p_0^{11} \cdot 0 = 0, \\
  \tau(p_0^4) &= 0 \cdot 0 - p_0^{11} \cdot (-p_0^{11}) = p_0^{22} \neq 0.
\end{align*}
Allgemein: $\tau(p_0^k) = 0$ für alle \emph{ungeraden} $k \geq 1$ und
$\tau(p_0^k) = (-1)^{k/2} p_0^{11k/2}$ für alle \emph{geraden} $k \geq 0$. Dieses
hochstrukturierte Verschwindensmuster wäre in der Theorie der Modulformen ohne
Präzedenzfall.

\subsection{Zusammenfassung der Konsequenzen}

\begin{proposition}\label{prop:konsequenzen}
  Wenn $\tau(p_0) = 0$ für eine Primzahl $p_0$ gilt, dann:
  \begin{enumerate}[label=(\roman*)]
    \item Der Eulerfaktor von $L(s,\Delta)$ bei $p_0$ hat Nullstellen auf der
          kritischen Geraden.
    \item $\tau(p_0^k) = 0$ für alle \emph{ungeraden} $k \geq 1$ und
          $\tau(p_0^k) = (-1)^{k/2} p_0^{11k/2}$ für alle geraden $k \geq 0$.
    \item Für alle Primzahlen $\ell \neq p_0$ ist $-p_0^{11}$ ein Quadrat modulo $\ell$
          (Serres Beobachtung).
    \item Die mod-$p_0$ Galois-Darstellung $\bar\rho_{p_0}$ ist reduzibel.
  \end{enumerate}
\end{proposition}

%% ============================================================
\section{Offene Probleme und Ausblick}
%% ============================================================

Die Lehmer-Vermutung ist eingebettet in ein Netz offener Probleme:

\begin{enumerate}[leftmargin=2em, label=\textbullet]
  \item \textbf{Lehmer-Vermutung selbst.} Beweis von $\tau(n) \neq 0$ für alle $n$.
        Kein unbedingter Beweis ist bekannt.
  \item \textbf{Nicht-Verschwinden von $L(k/2, f)$.} Gilt für eine allgemeine
        Hecke-Eigenform $f$ vom Gewicht $k$ stets $L(k/2, f) \neq 0$? Dies ist
        ein Analogon am Zentralpunkt.
  \item \textbf{Verteilung von $\tau(p)/2p^{11/2}$.} Die Sato-Tate-Vermutung
        (bewiesen von Taylor u.a.\ \cite{Taylor2008}) zeigt, dass diese Werte
        gleichverteilt auf $[-1, 1]$ gemäß dem Halbkreismaß sind. Das ist
        konsistent damit, dass $\tau(p) = 0$ ein Maß-null-Ereignis ist,
        schließt es aber nicht aus.
  \item \textbf{Effektive untere Schranke.} Kann man $|\tau(n)| \geq f(n)$ für
        eine explizite Funktion $f$ beweisen, die gegen unendlich strebt? Jede
        solche Schranke würde die Vermutung lösen.
  \item \textbf{Allgemeine Eigenformen.} Die Lehmer-Vermutung wird für alle
        Hecke-Eigenformen $f \in S_k(\Gamma)$ erwartet, nicht nur für $\Delta$;
        keiner dieser Fälle ist bewiesen.
\end{enumerate}

%% ============================================================
\section{Schlussfolgerung}
%% ============================================================

Die Lehmer-Vermutung $\tau(n) \neq 0$ für alle $n \geq 1$ ist eines der am
einfachsten formulierten, zugleich aber hartnäckigsten Probleme in der Arithmetik
der Modulformen. Wir haben den theoretischen Rahmen entwickelt, der sie stützt: die
Hecke-Eigenform-Struktur von $\Delta$, die Multiplikativität von $\tau$, das
Eulerprodukt für $L(s, \Delta)$, Delignes Satz über die Größe von $\tau(p)$ und die
strukturellen Konsequenzen, die eine Nullstelle nach sich ziehen würde. Unsere
numerische Verifikation im Specialist Mathematics Project bestätigt die Vermutung
bis $n = 50{.}000$ mit Sicherheit, und die beste bekannte Schranke liegt bei
$n \approx 10^{23}$.

Die Vermutung bleibt offen. Ihre Auflösung --- in welcher Richtung auch immer ---
würde einen bedeutenden Fortschritt in unserem Verständnis der Arithmetik von
Hecke-Eigenformen, mod-$\ell$ Galois-Darstellungen und des Nicht-Verschwindens
automorpher $L$-Funktionen darstellen.

%% ============================================================
\begin{thebibliography}{10}

\bibitem{Ramanujan1916}
S.~Ramanujan,
\textit{On certain arithmetical functions},
Trans.\ Cambridge Philos.\ Soc.\ \textbf{22} (1916), 159--184.

\bibitem{Mordell1917}
L.~J.~Mordell,
\textit{On Mr.\ Ramanujan's empirical expansions of modular functions},
Proc.\ Cambridge Philos.\ Soc.\ \textbf{19} (1917), 117--124.

\bibitem{Lehmer1947}
D.~H.~Lehmer,
\textit{The vanishing of Ramanujan's function $\tau(n)$},
Duke Math.\ J.\ \textbf{14} (1947), 429--433.

\bibitem{SwinnertonDyer1973}
H.~P.~F.~Swinnerton-Dyer,
\textit{On $\ell$-adic representations and congruences for coefficients of modular
forms},
in \textit{Modular Functions of One Variable III}, Lecture Notes in Math.\ \textbf{350},
Springer, 1973, S.\ 1--55.

\bibitem{Deligne1974}
P.~Deligne,
\textit{La conjecture de Weil. I},
Publ.\ Math.\ Inst.\ Hautes \'Etudes Sci.\ \textbf{43} (1974), 273--307.

\bibitem{DelignePureI}
P.~Deligne,
\textit{La conjecture de Weil. II},
Publ.\ Math.\ Inst.\ Hautes \'Etudes Sci.\ \textbf{52} (1980), 137--252.

\bibitem{Serre1987}
J.-P.~Serre,
\textit{Quelques applications du th\'eor\`eme de densit\'e de Chebotarev},
Publ.\ Math.\ Inst.\ Hautes \'Etudes Sci.\ \textbf{54} (1981), 123--201.

\bibitem{Chua2004}
K.~S.~Chua,
\textit{Extremal modular forms and the Lehmer conjecture},
Preprint, 2004.

\bibitem{Taylor2008}
R.~Taylor,
\textit{Automorphy for some $l$-adic lifts of automorphic mod $l$ Galois
representations. II},
Publ.\ Math.\ Inst.\ Hautes \'Etudes Sci.\ \textbf{108} (2008), 183--239.

\bibitem{DiamondShurman}
F.~Diamond und J.~Shurman,
\textit{A First Course in Modular Forms},
Graduate Texts in Mathematics \textbf{228}, Springer, 2005.

\end{thebibliography}

\end{document}
